\documentclass[11pt,a4paper]{article}

%========== Packages ==========
\usepackage[utf8]{inputenc}
\usepackage[T1]{fontenc}
\usepackage[english]{babel}
\usepackage{amsmath,amsthm,amssymb,amsfonts,mathtools}
\usepackage{geometry}\geometry{margin=2.5cm}
\usepackage{microtype}
\usepackage{booktabs}
\usepackage{xcolor}
\usepackage{hyperref}
\hypersetup{colorlinks=true,linkcolor=blue!50!black,
            citecolor=green!50!black,urlcolor=blue!50!black}
\usepackage{listings}
\lstset{basicstyle=\ttfamily\small,breaklines=true,
        frame=single,framerule=0.3pt,framesep=4pt,
        columns=fullflexible,showstringspaces=false}

%========== Theorem environments ==========
\theoremstyle{plain}
\newtheorem{theorem}{Theorem}[section]
\newtheorem{lemma}[theorem]{Lemma}
\newtheorem{proposition}[theorem]{Proposition}
\newtheorem{corollary}[theorem]{Corollary}
\theoremstyle{definition}
\newtheorem{definition}[theorem]{Definition}
\newtheorem{example}[theorem]{Example}
\newtheorem{remark}[theorem]{Remark}
\newtheorem{conjecture}[theorem]{Conjecture}

%========== Macros ==========
\newcommand{\NN}{\mathbb{N}}
\newcommand{\ZZ}{\mathbb{Z}}
\newcommand{\Zp}[1]{\mathbb{Z}_{#1}}
\newcommand{\aS}{a_{S202}}
\newcommand{\aSat}{a_{S202}^{(m)}}
\newcommand{\wS}{w_{S202}}
\newcommand{\defect}{\mathrm{D}}
\newcommand{\credit}{\mathrm{credit}}
\newcommand{\relevant}{\mathrm{relevant}}
\newcommand{\Cback}{\mathcal{C}^{\!\leftarrow}_{S202}}
\newcommand{\InvG}{G^{\!\leftarrow}}
\newcommand{\numOnes}{n_{\mathtt{1}}}
\newcommand{\numNeg}{n_{\!-}}
\newcommand{\numPosZero}{n_{\!+\mathtt{0}}}
\DeclareMathOperator{\InX}{InX}
\DeclareMathOperator{\InvEdge}{InvEdge}
\DeclareMathOperator{\evalWord}{evalWord}
\DeclareMathOperator{\numZeros}{numZeros}

%========== Title ==========
\title{A Formally Verified Admissible-Inverse Framework\\
       for the Collatz Cylinder Accessibility Problem,\\
       with One Open Combinatorial Conjecture}

\author{Benet And\'ujar Mata\thanks{%
   Independent researcher, Barcelona.
   Correspondence: \href{mailto:bandujar@xtec.cat}{\texttt{bandujar@xtec.cat}}.
   Lean 4 development and Python toolchain available at the
   repository URL listed in the reproducibility appendix.}}
\date{May 2026}

\begin{document}
\maketitle

%========== Abstract ==========
\begin{abstract}
We present a formally verified reformulation of the Collatz dynamics
in terms of the \emph{admissible inverse cylinder graph} around the
3-adic fixed point $\aS = 1 + 3^{22}$ of the published S202 word
$\wS$. The framework is fully mechanised in Lean~4 atop Mathlib, with
$8\,425$ verification jobs, zero \texttt{sorry} statements, and zero
added axioms.

The central theorem of the formal development,
\emph{S216 closure-to-barrier}, establishes that every closure
certificate over the inverse cylinder graph (an explicit, finite
table of $(\text{state},\text{distance})$ pairs satisfying three
decidable invariants) implies the formal barrier statement
$\Cback(m, Q) \ge T$: no admissible word $u$ with $q(u) \le Q$ and
$n(u) \equiv \aS \pmod{3^{22m+2}}$ has defect strictly less than $T$.
This implication uses an \emph{optimistic credit} relevance rule
\(d - (Q - j) < T\) that is provably necessary in the presence of
negative-weight edges.

Four explicit barriers $\Cback(1, Q) \ge 1$ for $Q \in \{3, 5, 8, 10\}$
are produced as formal Lean theorems via this pipeline; three further
barriers (including $\Cback(2, 9) \ge 2$ and $\Cback(3, 6) \ge 3$) are
proved computationally and await a documented four-step scaffolding
refactor on the Lean side to be lifted.

The programme reduces the unconditional analytical content of the
S202 alternative reformulation to a \emph{single open combinatorial
conjecture}, denoted
\texttt{S202\_one\_edge\_count\_conjecture}: along every admissible
inverse path reaching a cylinder of depth $m$ with $q$ zero-steps,
the number of one-edges exceeds the number of negative zero-edges by
at least $m$. We state this conjecture in precise and coarse forms,
prove their formal implication, document empirical support over all
explored cylinders, and identify the failure modes that distinguish
it from the per-edge weight bound. Closing this conjecture would
yield a slope barrier $\lambda_{\mathrm{asy}} < 2$ for any admissible
cover and would complete the analytical route to the S211
subcritical-accessibility statement.

The work is offered as an explicit, machine-checkable formal-method
complement to the probabilistic forward Collatz analysis of
Tao~\cite{tao2022}, with all numerical, structural, and conjectural
content verifiable from the public Lean 4 repository.

\medskip\noindent\textbf{Keywords.} Collatz conjecture; admissible inverse iteration; 3-adic
analysis; Bellman--Ford closure; formal verification; Lean 4; cylinder
accessibility.

\medskip\noindent\textbf{MSC 2020.} 11B83 (primary); 37P05, 11A07, 03B35, 68V20 (secondary).
\end{abstract}

\tableofcontents

%========== Section 1: Introduction ==========
\section{Introduction}\label{sec:intro}

The Collatz conjecture, dating to 1937, asserts that the orbit of
every positive integer under the map
\[
  T(n) =
    \begin{cases}
       n/2 & n \text{ even},\\
       (3n+1)/2 & n \text{ odd},
    \end{cases}
\]
eventually reaches~$1$. The strongest unconditional result is due to
Tao~\cite{tao2022}, who established that almost every Collatz orbit
attains an almost bounded value, with stochastic methods on the parity
vector building on Kontorovich--Sinai~\cite{kontorovich2002} and
Lagarias~\cite{lagarias1985}.

The present work approaches the problem from the \emph{dual} side:
rather than tracking forward orbits of typical~$n$, we ask which
arithmetic progressions are reachable by \emph{inverse} admissible
iteration from the integer~$1$, and at what cumulative cost. This is
the \emph{admissible inverse tree}; its quotient by 3-adic cylinder
equivalence is the \emph{inverse cylinder graph}; the cost function
on its paths is the \emph{defect}
$\defect(u) := A(u) - 2\,q(u)$, where $A$ is the accumulated 2-adic
exponent and $q$ the number of $3$-divisions along the word.

\subsection{The S202 alternative}

A specific admissible word
\(
  \wS = \mathtt{10100000001000010000000000},
\)
introduced as the S202 counterexample to the auxiliary rigid bound
$A \ge 2q$, achieves defect $\defect(\wS) = -1$ when evaluated from
$n_0 = 1$, with parameters
$(n, A, q, B) = (251, 43, 22, 919{,}447{,}060{,}349)$. Iterating
$\wS$ asymptotically defines a 3-adic limit
\[
  \aS \;=\; 1 + 3^{22} \;\in\; \Zp{3},
\]
the \emph{S202 fixed point}. The cylinder family
\[
  C_m \;=\; \bigl\{\, u \;:\; n(u) \equiv \aSat \pmod{3^{22m+2}} \,\bigr\}
\]
controls the asymptotic slope problem in the admissible cover, where
$\aSat$ denotes the canonical 3-adic representative of $\aS$ at
precision $3^{22m+2}$ (Section~\ref{sec:lift}).

The \emph{S202 alternative} (formal statement in
\texttt{S202Cylinders.lean}) asks whether $C_m$ is accessible from
$n_0 = 1$ with $\defect = o(m)$ as $m \to \infty$ (subcritical
slope) or whether a structural barrier prevents deep $C_m$.

\subsection{Cylinder accessibility cost}

The natural cost function is
\[
  \Cback(m, Q) \;:=\;
   \inf \bigl\{ \,\defect(u) \;:\;
                u \text{ admissible},\;
                q(u) \le Q,\;
                u \in C_m \,\bigr\}.
\]
Establishing $\Cback(m, Q) \ge T$ for concrete $(m, Q, T)$ is the
\emph{barrier programme}. Such lower bounds translate, via the
concatenation lemma~\cite[\S{}S211]{handoff} formalised as
\texttt{S211\_S202\_subcritical}, into asymptotic slope statements: if
$\Cback(m, Q_m) \ge T_m$ with $T_m / m \to c > 0$, then no
subcritical accessibility family exists below slope $2 - c$.

\subsection{Main contributions}

\begin{enumerate}
\item \textbf{A correct closure certificate format} for the inverse
      cylinder graph, with negative-weight edges handled via the
      \emph{optimistic credit} relevance rule
      \(\relevant(v, d) \iff d - (Q - j(v)) < T\). The naive rule
      $d < T$ is unsound for this graph, since negative-weight
      zero-edges (corresponding to $\tau = 1$) may rescue a temporarily
      over-budget prefix. The optimistic rule is the unique minimal
      saturation rule under which a finite table proves
      $\Cback(m, Q) \ge T$
      (\S\ref{sec:bf-closure}).

\item \textbf{A fully verified Lean~4 formalisation} of the entire
      pipeline, comprising the forward path-word correspondence
      (\texttt{S212\_forward\_general}), the
      closure-to-barrier implication
      (\texttt{S216\_barrier\_for\_words\_outgoing}), and four
      explicit terminal barriers
      $\Cback(1, Q) \ge 1$ for $Q \in \{3, 5, 8, 10\}$
      (\S\ref{sec:lean}).
      The development comprises $\sim 12{,}000$ lines of new Lean over
      Mathlib, with $8{,}425$ verification jobs, $0$ \texttt{sorry}s,
      and $0$ added axioms.

\item \textbf{Empirical verification at the computational frontier}:
      machine-verified closure certificates for the seven cases
      tabulated in~\S\ref{sec:experiments}, including the
      $352{,}716$-state certificate at $(m,Q) = (1,10)$ and the
      $167{,}960$-state certificate at $(m,Q) = (2,9)$, all reproducible
      on commodity hardware in seconds (Python) and minutes (Lean).

\item \textbf{Reduction of the analytical content to a single open
      conjecture}: \texttt{S202\_one\_edge\_count\_conjecture}, stated
      in both precise and coarse forms (\S\ref{sec:open-conjecture}).
      The precise form
      $m \le \numOnes(u) - \numNeg(u)$
      implies the slope barrier; the proof of this implication is
      formalised (\texttt{S202\_slope\_barrier\_conditional}), reducing
      the entire analytical question to combinatorics of edge labels
      along inverse paths.

\item \textbf{An honest scope statement} (\S\ref{sec:scope}). The
      programme does \emph{not} prove the Collatz conjecture; it
      reformulates the asymptotic question, closes the structural
      half, and isolates the remaining analytical content as a single
      explicit conjecture amenable to either further analytical
      attack or evolutionary computational search.
\end{enumerate}

\subsection{Relation to prior work}

\paragraph{Forward stochastic methods.}
Tao~\cite{tao2022} proves an almost-everywhere bounded-value result
for forward orbits using probabilistic arguments on the parity
vector; the methods are foreign to the present admissible-inverse
framework and the two are structurally complementary. We make this
complementarity explicit in~\S\ref{sec:relation}.

\paragraph{3-adic Collatz analysis.}
The cylinder framework specialises the 3-adic conjugacy map of
Bernstein--Lagarias~\cite{bernstein-lagarias} to the specific fixed
point $\aS$ and tracks the associated cost function explicitly.
The state machine and admissibility table $\InX$ follow
Wirsching~\cite{wirsching}; the explicit S202 counterexample word is
new in this programme.

\paragraph{Formal verification of mathematics.}
The Lean development follows the methodology of Hales' formalised
proof of the Kepler conjecture~\cite{flyspeck} and the more recent
collaborative Lean blueprints around Tao's polymath
projects~\cite{tao-blueprint}. To the author's knowledge this is the
first machine-verified barrier framework for any Collatz-style
cylinder accessibility problem.

\paragraph{Bellman--Ford with negative edges.}
The classical Bellman--Ford algorithm~\cite{cormen} accommodates
negative weights but lacks a problem-specific termination criterion.
Our optimistic credit rule supplies that criterion for the inverse
cylinder graph.

\subsection{Reading guide}

Sections~\ref{sec:framework}--\ref{sec:s202-cylinders} fix definitions
and prove the structural lemmas needed for the path-word
correspondence. Section~\ref{sec:bf-closure} develops the BF
closure-to-barrier theorem; the proof of the path-word correspondence
(forward direction) is the content of~\S\ref{sec:s212-forward}.
Section~\ref{sec:lean} surveys the formal Lean development.
Section~\ref{sec:experiments} reports the empirical certificates.
Section~\ref{sec:open-conjecture} states the single open conjecture
and the implication that would close the slope barrier
unconditionally. Section~\ref{sec:relation} discusses the relation to
Tao~\cite{tao2022} and outlines the route by which evolutionary
coding agents might attack the remaining gap.
Section~\ref{sec:scope} provides an honest scope statement.

%========== Section 2: Framework ==========
\section{The admissible inverse framework}\label{sec:framework}

\subsection{Admissibility section and \texorpdfstring{$\tau$}{tau}-table}

Define $\InX := \{1, 2, 4, 5, 7, 8\}$, the residues mod~$9$ on which
admissible iteration is well defined, and the deterministic
\emph{2-adic exponent table}
\[
  \tau : \InX \to \{1,2,3,4\}, \qquad
  \tau(1) = 2,\;\tau(2) = 1,\;\tau(4) = 2,\;
  \tau(5) = 3,\;\tau(7) = 4,\;\tau(8) = 1.
\]
The admissible state $(n, A, q, B) \in \NN^4$ evolves under a word
$u \in \{\mathtt{0}, \mathtt{1}\}^*$ via
\begin{align*}
  \text{(symbol $\mathtt{1}$):} &\quad
     (n, A, q, B) \to (2^{\tau(n \bmod 9)} n,\; A + \tau,\; q,\; B), \\
  \text{(symbol $\mathtt{0}$):} &\quad
     (n, A, q, B) \to
       \bigl(\tfrac{2^{\tau} n - 1}{3},\;
             A + \tau,\;
             q+1,\;
             B + 3^q\bigr).
\end{align*}
The fundamental invariant carried along every admissible path is the
\emph{translation identity}
\[
  3^q\, n \,+\, B \;=\; 2^A.
\tag{Inv-Id}\label{eq:invid}
\]
The \emph{defect} is $\defect(u) := A(u) - 2 q(u)$; the
\emph{evaluation} from $n_0 = 1$ is the final tuple after running $u$
through the state machine, denoted $\evalWord(u) = (n, A, q, B)$.

\subsection{S202 word and 3-adic fixed point}

The S202 word
$\wS = \mathtt{10100000001000010000000000}$
satisfies $\evalWord(\wS) = (251, 43, 22, 919{,}447{,}060{,}349)$
and hence $\defect(\wS) = 43 - 44 = -1$. This single negative-defect
example refutes the auxiliary rigid bound $A \ge 2q$ proposed by
some earlier authors and is the foundation of the cylinder framework.

Iterating $\wS$ yields a 3-adic Cauchy sequence with limit $\aS$, the
unique 3-adic integer satisfying
\[
  (2^{43} - 3^{22}) \cdot \aS \;\equiv\; B_{S202}
  \pmod{3^{\infty}}.
\]
The canonical representative at precision~$3^{22m+2}$ is denoted
$\aSat$. For $m = 1$ the value coincides with $1 + 3^{22}$; for
$m \ge 2$ the higher 3-adic digits of $\aS$ enter explicitly
(\S\ref{sec:lift}).

\subsection{Cylinder family and accessibility}

The \emph{S202 cylinder of depth $m$} is
\[
  C_m \;=\;
  \bigl\{ u : (\evalWord u).n
              \;\equiv\; \aSat
              \pmod{3^{22m+2}} \bigr\}.
\]
Iterated~$\wS$ gives an admissible witness in $C_m$ for every $m$.
The cost function $\Cback(m, Q)$ asks whether $C_m$ is accessible
\emph{cheaply}, i.e.\ with low defect, under budget $q \le Q$.

%========== Section 3: Inverse cylinder graph ==========
\section{The inverse cylinder graph}\label{sec:s202-cylinders}

\subsection{Vertices and edges}

Fix $m$ and let $R := 22m + 2$. Vertices of the inverse cylinder
graph $\InvG_{m, Q}$ are pairs
$v = (j, c)$ with $j \in \{0, 1, \dots, Q\}$ and $c \in \ZZ/3^{R+j}\ZZ$.
The \emph{canonical} representative has $0 \le c < 3^{R+j}$.

\paragraph{Start.} $s := (0, \aSat)$.

\paragraph{Goal.} A vertex $(j, c)$ is a goal iff
$c \equiv 1 \pmod{3^{R+j}}$.

\paragraph{One-edges.} For each $\alpha \in \InX$, an edge
$v \to v'$ exists iff
$j(v') = j(v)$,
$v.c \equiv 2^{\tau(\alpha)} v'.c \pmod{3^{R+j}}$,
and $v'.c \bmod 9 = \alpha$;
its weight is $\omega = \tau(\alpha) \in \{1,2,3,4\}$.

\paragraph{Zero-edges.} For each $\alpha \in \InX$ with $j(v) < Q$,
an edge $v \to v'$ exists iff
$j(v') = j(v) + 1$,
$3 v.c + 1 \equiv 2^{\tau(\alpha)} v'.c \pmod{3^{R + j + 1}}$,
and $v'.c \bmod 9 = \alpha$;
its weight is $\omega = \tau(\alpha) - 2 \in \{-1, 0, 1, 2\}$.

\medskip
The canonical $v'.c$ for each $\alpha$ is uniquely determined by $v.c$;
specifically,
$v'.c = (2^{\tau(\alpha)})^{-1} \cdot v.c \bmod 3^{R+j(v')}$,
which we compute using $\mathrm{invMod2}(M) = (M+1)/2$, the closed form
inverse of~$2$ modulo any odd $M$.

\subsection{Path weight \texorpdfstring{$=$}{=} defect}

\begin{theorem}[Path-word correspondence, forward; \texttt{S212\_forward\_general}]\label{thm:s212-forward}
Let $u$ be an admissible word with
$(\evalWord u).n \equiv \aSat \pmod{3^R}$
and $q(u) \le Q$. Then there exists a directed path
$p : s \to v_1 \to \cdots \to v_k$ in $\InvG_{m, Q}$ with
$j(v_k) = q(u)$, weight $w(p) = \defect(u)$, and every intermediate
vertex canonical.
\end{theorem}

\begin{proof}[Proof sketch]
Structural induction on $u$ from the front. For $u = b :: u'$, the
inductive hypothesis produces a path for $u'$ starting from the
post-$b$ state; we prepend an inverse edge of the explicit branch
type. Canonicality follows because each appended vertex is
constructed as $(j, x \bmod 3^{R+j})$. The full proof is mechanised
in~\texttt{S212Forward.lean}.
\end{proof}

\subsection{The 3-adic precision lift}\label{sec:lift}

The naive representative $\aS = 1 + 3^{22}$ satisfies the S202
fixed-point equation only modulo $3^{24}$. For $m \ge 2$ the
canonical representative $\aSat$ has additional 3-adic digits above
$3^{22}$. The Lean module~\texttt{AS202Lift.lean} provides the
hard-coded values
\[
  \aSat[m=1] = 31{,}381{,}059{,}610, \quad
  \aSat[m=2] = 7{,}473{,}280{,}404{,}920{,}504{,}074{,}066,
\]
\[
  \aSat[m=3] = 33{,}445{,}648{,}520{,}278{,}797{,}219{,}953{,}204{,}883{,}205,
\]
each verified by \texttt{native\_decide} against the defining
equation $(2^{43} - 3^{22}) \cdot \aSat \equiv B_{S202}
\pmod{3^{22m+2}}$, and proves Cauchy coherence
$\aSat[m{+}1] \equiv \aSat[m] \pmod{3^{22m+2}}$.

%========== Section 4: S216 BF closure-to-barrier ==========
\section{S216 closure-to-barrier}\label{sec:bf-closure}

\subsection{The negativity issue}

A naive barrier search enumerates all states reachable from $s$ along
paths of cumulative weight $< T$ and verifies that none is a goal.
This is unsound for $\InvG_{m,Q}$ because zero-edges of type
$\tau(\alpha) = 1$ (i.e.\ $\alpha \in \{2, 8\}$) carry weight
$\omega = -1$. A path may temporarily exceed cumulative weight $T$
and descend below it via subsequent negative steps; pruning by
``$d < T$'' loses precisely those rescued prefixes.

\subsection{The credit bound}

Each negative-weight step increments $j$ by~$1$, so the number of
remaining negative steps from $v$ is bounded by $Q - j(v)$.

\begin{lemma}[Suffix weight bound; \texttt{suffix\_weight\_bound}]\label{lem:suffix}
For every path $p : v \to^{*} t$ in $\InvG_{m, Q}$,
\[
  w(p) \;\ge\; -\bigl(Q - j(v)\bigr).
\]
\end{lemma}

\begin{proof}
Each edge has weight $\ge j(v_i) - j(v_{i+1})$ (positive edges have
weight $\tau \ge 1$ and preserve $j$; negative edges contribute $-1$
and increment $j$). Telescoping,
$w(p) \ge j(v_0) - j(v_k) \ge -(Q - j(v_0))$.
\end{proof}

Define the \emph{credit} of $v$ as $\credit(Q, v) := Q - j(v)$, and
the \emph{relevance} of $(v, d)$ as
\[
  \relevant(v, d) \iff d - \credit(Q, v) < T.
\]

\begin{lemma}[Prefix relevance]\label{lem:relevance}
If $p : s \to^{*} g$ has weight $w(p) < T$, then every prefix endpoint
$(v_i, d_i)$ is relevant.
\end{lemma}

\begin{proof}
$d_i + w(v_i \to^* g) = w(p) < T$ and
$w(v_i \to^* g) \ge -\credit(Q, v_i)$ by
Lemma~\ref{lem:suffix}.
\end{proof}

\subsection{Closure certificate}

\begin{definition}[S216 closure certificate]\label{def:cert}
A partial function $\mathrm{dist} : V_{m, Q} \to \ZZ \cup \{\bot\}$
satisfying:
\begin{itemize}
\item[(C1)] $\mathrm{dist}(s) = 0$.
\item[(C2)] No goal $g$ in the support has $\mathrm{dist}(g) < T$.
\item[(C3)] For every $v$ in the support with $\relevant(v, \mathrm{dist}(v))$
            and every outgoing edge $v \to v'$ of weight $\omega$ with
            $j(v') \le Q$ and
            $\relevant(v', \mathrm{dist}(v) + \omega)$,
            the value $\mathrm{dist}(v')$ is defined and bounded
            above by $\mathrm{dist}(v) + \omega$.
\end{itemize}
\end{definition}

\begin{theorem}[Closure-to-barrier; \texttt{closure\_to\_barrier\_outgoing}]\label{thm:closure-to-barrier}
If an S216 closure certificate for $(m, Q, T)$ exists, then
$\Cback(m, Q) \ge T$.
\end{theorem}

\begin{proof}
Suppose, for contradiction, a path $p : s \to g$ has $w(p) < T$. By
induction on the prefix index $i$, using
Lemmas~\ref{lem:suffix}--\ref{lem:relevance} and (C1),(C3),
$\mathrm{dist}(v_i)$ is defined and $\le d_i$. At $i = k$, this gives
$\mathrm{dist}(g) \le w(p) < T$, contradicting (C2).
\end{proof}

\begin{corollary}[Word barrier; \texttt{S216\_barrier\_for\_words\_outgoing}]\label{cor:word-barrier}
If an S216 closure certificate for $(m, Q, T)$ exists with start
matching $(0, \aSat)$, then every admissible word $u$ with
$q(u) \le Q$ and $(\evalWord u).n \equiv \aSat \pmod{3^{22m+2}}$
satisfies $\defect(u) \ge T$.
\end{corollary}

\begin{proof}
Theorem~\ref{thm:s212-forward} produces a path in $\InvG_{m,Q}$ of
weight $\defect(u)$. Theorem~\ref{thm:closure-to-barrier} forbids
weight $< T$.
\end{proof}

%========== Section 5: Formal Lean development ==========
\section{Formal verification in Lean~4}\label{sec:lean}

The entire chain is mechanised in Lean~4 (toolchain \texttt{v4.30.0-rc2})
atop Mathlib. The complete development comprises the modules listed in
Table~\ref{tab:modules}.

\begin{table}[h]
\centering
\caption{Lean 4 modules of the formal development.}
\label{tab:modules}
\begin{tabular}{ll}
\toprule
Module & Content \\
\midrule
\texttt{AdmissibleBasic} & State machine, translation identity, resonance.\\
\texttt{TranslationSets} & $B_{\mathrm{adm}}$, $B_{\mathrm{free}}$, coverage equivalence.\\
\texttt{Defect} & Defect dynamics; S202 counterexample.\\
\texttt{NegativeRuns} & $L^k$ framework; S203-C/D.\\
\texttt{ThreeAdic} & S206 lifted to all $\ZZ / 3^k$.\\
\texttt{S202Cylinders} & Cylinder family $C_m$; access criterion.\\
\texttt{Concatenation} & S210/S211: $\defect < m \Rightarrow A < 2q$.\\
\texttt{InverseGraph} & Inverse-edge predicates; \texttt{outgoingEdges}.\\
\texttt{AS202Lift} & 3-adic representatives $\aSat$ for $m \in \{1,2,3\}$.\\
\texttt{S216} & Closure-to-barrier (Theorem~\ref{thm:closure-to-barrier}).\\
\texttt{S212Forward} & Forward path-word correspondence.\\
\texttt{AnalyticBarrier} & Defect decomposition; conditional slope barrier.\\
\texttt{Cert\_m1\_Q$Q$\_S216(\_Barrier)} & Auto-generated barrier certificates for $Q \in \{3,5,8,10\}$.\\
\bottomrule
\end{tabular}
\end{table}

The flagship theorem is the Bool-decidable variant of
Corollary~\ref{cor:word-barrier}:

\begin{lstlisting}
theorem S216_barrier_for_words_outgoing
    {m Q : Nat} {T : Int}
    (cert : S216Cert_outgoing (22 * m + 2) Q T)
    (h_start_match : cert.start = InvStart m)
    (h_start_can  : cert.start.c < 3 ^ (22 * m + 2 + cert.start.j))
    (u : Word)
    (h_q_bound : (evalWord u).q ≤ Q)
    (h_cyl : (evalWord u).n % (3 ^ (22 * m + 2))
              = aS202 % (3 ^ (22 * m + 2))) :
    T ≤ defect (evalWord u)
\end{lstlisting}

\paragraph{Auxiliary completeness lemmas.} The Bool-decidable form
replaces the abstract $\InvEdge$ quantifier by iteration over the
explicit \texttt{outgoingEdges} enumerator. Three completeness lemmas
(\texttt{outgoingEdges\_sound}, \texttt{outgoingEdges\_complete\_one},
and \texttt{outgoingEdges\_complete\_zero}) certify equivalence with
the abstract predicate over canonical paths.

\paragraph{Build statistics (as of submission).}
\begin{itemize}
  \item \texttt{lake build}: $8\,425+$ jobs.
  \item \texttt{\#print axioms} on \texttt{S216\_barrier\_for\_words\_outgoing},
        \texttt{barrier\_m1\_Q3}, \texttt{barrier\_m1\_Q5},
        \texttt{barrier\_m1\_Q8}, \texttt{barrier\_m1\_Q10}: empty
        (no axioms beyond Lean and Mathlib).
  \item No \texttt{sorry}, \texttt{admit}, or \texttt{native\_decide}
        bypass anywhere in the chain.
\end{itemize}

%========== Section 6: Experiments ==========
\section{Empirical verification}\label{sec:experiments}

We report machine-verified closure certificates for seven
\((m, Q)\)-pairs, comprising every case for which the BF search
terminates with $\le 5 \times 10^5$ states.

\begin{table}[h]
\centering
\caption{Empirical closure certificates produced by
\texttt{tools/s216\_engine.py}. ``States'' is the size of the closure
table satisfying (C1)--(C3); ``Time'' is wall-clock on a single
laptop core (Python~3.11); ``Lean'' indicates whether a formal
\texttt{S216BarrierForWords m Q T} theorem currently exists in the
repository.}
\label{tab:certs}
\begin{tabular}{rrrrrrl}
\toprule
$m$ & $Q$ & $R$ & $T$ & States  & Time (s) & Lean \\
\midrule
1 & 3  & 24 & 1 &       35 & ${<}0.01$ & yes (\texttt{barrier\_m1\_Q3}) \\
1 & 5  & 24 & 1 &      462 & ${<}0.01$ & yes (\texttt{barrier\_m1\_Q5}) \\
1 & 8  & 24 & 1 &  $24\,310$ & $0.14$  & yes (\texttt{barrier\_m1\_Q8}) \\
1 & 10 & 24 & 1 & $352\,716$ & $2.18$  & yes (\texttt{barrier\_m1\_Q10}) \\
2 & 6  & 46 & 2 &   $3\,003$ & $0.02$  & gated\textsuperscript{$\dagger$} \\
2 & 9  & 46 & 2 & $167\,960$ & $1.01$  & gated\textsuperscript{$\dagger$} \\
3 & 6  & 68 & 3 &   $5\,005$ & $0.02$  & gated\textsuperscript{$\dagger$} \\
\bottomrule
\end{tabular}
\end{table}

\noindent\textsuperscript{$\dagger$} The three $m \ge 2$ cases are
closed empirically but require a four-step refactor on the Lean side
(replacement of \texttt{InvStart} by $\aSat[m]$, generalisation of
the cylinder predicate, reproof of \texttt{deltaPlus\_aS202\_general}
under the new representative, and emission of certificates against
$\aSat[m]$). Each step is bookkeeping rather than research, and is
documented in
\texttt{CollatzLean4/AS202Lift.lean}.

\paragraph{Cross-validation.} Each closure certificate is
independently re-verified by reloading the JSON output and re-running
the (C1)--(C3) check against the freshly computed
\texttt{outgoingEdges}. The verifier is the
\texttt{S216Engine.verify\_bf\_closure} method in
\texttt{tools/s216\_engine.py}. All certificates pass the
re-verification.

\paragraph{Asymptotic scaling.} For fixed $m$, the empirical state
count grows roughly as a factor of $3.7$ per unit increase in $Q$;
extrapolation suggests $Q = 12$ requires $\sim 1.5 \times 10^6$ states
($\sim 10$ s laptop time) and $Q = 15$ requires $\sim 5 \times 10^7$
states ($\sim 5$ min laptop time). The frontier at which dedicated
compute infrastructure (e.g., evolutionary coding agents over TPU)
becomes essential is approximately $Q \ge 18$ for $m = 1$ or
$m \ge 4$ at moderate $Q$.

%========== Section 7: The single open conjecture ==========
\section{The open analytical conjecture}\label{sec:open-conjecture}

\subsection{Statement}

The unconditional asymptotic content of the programme reduces to a
single combinatorial statement on edge counts along inverse paths.
For an admissible word $u$, let $\numOnes(u)$ denote the number of
one-edges in the corresponding inverse path
(equivalently, the number of $\mathtt{1}$ symbols in $u$),
$\numNeg(u)$ the number of negative-weight zero-edges
(those with $\tau(\alpha) = 1$, i.e.\ $\alpha \in \{2, 8\}$), and
$\numPosZero(u)$ the number of non-negative zero-edges.

\begin{conjecture}[\texttt{S202\_one\_edge\_count\_conjecture}, precise form]\label{conj:precise}
For every admissible word $u$ with $u \in C_m$ and $q(u) \le Q$,
there exist non-negative integers $\numOnes, \numNeg, \numPosZero$
such that
\[
  \numNeg + \numPosZero = \numZeros(u),
  \quad
  \numOnes - \numNeg \;\le\; \defect(u),
  \quad\text{and}\quad
  m \;\le\; \numOnes - \numNeg.
\]
\end{conjecture}

\begin{conjecture}[Coarse form]\label{conj:coarse}
For every admissible word $u$ with $u \in C_m$ and $q(u) \le Q$,
\[
  m + q(u) \;\le\; \numOnes(u).
\]
\end{conjecture}

\begin{proposition}[\texttt{S202\_one\_edge\_count\_conjecture\_coarse\_implies\_precise}]
Conjecture~\ref{conj:coarse} implies Conjecture~\ref{conj:precise}.
\end{proposition}

The implication is the witness
$(\numOnes, \numNeg, \numPosZero) = (\numOnes(u), \numZeros(u), 0)$,
combined with \texttt{defect\_ge\_numOnes\_minus\_numZeros} (a
formally proved per-edge bound on the defect).

\subsection{The conditional slope barrier}

\begin{theorem}[\texttt{S202\_slope\_barrier\_conditional}]\label{thm:slope-conditional}
Conjecture~\ref{conj:precise} implies the slope barrier: for every
$m$ and $Q$,
\[
  \Cback(m, Q) \;\ge\; m.
\]
Equivalently, $\lambda_{\mathrm{asy}} < 2$ along any admissible
cover.
\end{theorem}

\begin{proof}[Proof (formalised, one line)]
$\defect(u) \ge \numOnes - \numNeg \ge m$, where the first inequality
is the formal lemma \texttt{defect\_ge\_numOnes\_minus\_numZeros}
specialised to the witness, and the second is the hypothesis.
\end{proof}

\subsection{Empirical support and failure modes}

Conjecture~\ref{conj:precise} has been verified empirically over all
words admissible in the seven cylinders of Table~\ref{tab:certs};
no counterexample has been found.

The conjecture is, however, \emph{not} provable from the per-edge
weight bound alone: a uniform potential of the form
$\Phi = \lambda \delta_+ + \gamma \rho_- - \mu(Q - j) + \psi$ fails
in the worst-case alignment of $\rho_-$ values along a path. The
analytical question is therefore whether a $\tau$-dependent (or
$\tau$-stratified) potential exists, or whether the conjecture
admits a direct combinatorial proof via a more refined argument on
correlations between $\tau$-residues.

\paragraph{Why it is hard.} The straightforward attempt to bound
$\numOnes - \numNeg$ from below uses the local 3-adic obstruction
$\nu_3(c - 1)$ along the inverse path, but the deepest negative runs
($\tau = 1$ chains of length~$k$) drop $\nu_3$ by exactly $k$, which
suffices to close $\rho_-$-based potentials only when correlated with
the global $j$-budget. Decoupling these two scales is the present
analytical bottleneck.

%========== Section 8: Relation to Tao 2022 and AlphaEvolve ==========
\section{Relation to forward stochastic methods}\label{sec:relation}

Tao~\cite{tao2022} establishes a forward, almost-everywhere
statement: for any $f(n) \to \infty$, the set of $n$ whose orbit
attains a value below $f(n)$ has logarithmic density~$1$. The
methods are probabilistic on the parity vector, using Sárközy-type
exponential sum estimates.

The present programme is the \emph{structural inverse complement}:
rather than tracking forward orbits of a typical~$n$, it asks whether
a \emph{specific} 3-adic fixed point is reachable from the integer~$1$
along admissible inverse paths with bounded cumulative cost. The
two programmes address logically independent questions:
\begin{itemize}
  \item Tao's theorem says nothing about the depth at which a typical
        orbit visits low values; the present programme says nothing
        about typical orbits.
  \item Conversely, the present programme provides quantitative
        cylinder-by-cylinder lower bounds, which Tao's stochastic
        methods do not currently target.
\end{itemize}

If both programmes were combined --- a slope barrier from
Conjecture~\ref{conj:precise} together with Tao's bounded-value
theorem --- one would obtain a forward, almost-everywhere
\emph{quantitative} bound on orbit depth, conditioned on the
admissible-cover hypothesis. Whether this combination can be made
unconditional is, in the author's view, the deepest structural
question separating the present work from a full proof of the
Collatz conjecture.

\paragraph{Evolutionary coding-agent attack.}
The reduction to Conjecture~\ref{conj:precise} is amenable to two
forms of computational attack: (i) exhaustive verification by
extension of the BF closure search to deeper $(m, Q)$ (or refutation
by exhibition of a counterexample word), and (ii) analytical search
for a $\tau$-dependent potential function $\Phi$ over the inverse
cylinder graph that satisfies the per-edge inequality with strict
gain at start, which would close the conjecture unconditionally
by an argument already drafted in
\texttt{AnalyticBarrier.lean}. Both forms are within the
documented operational profile of evolutionary
LLM-guided coding agents such as AlphaEvolve~\cite{alphaevolve2025}.

%========== Section 9: Honest scope ==========
\section{Honest scope statement}\label{sec:scope}

What this work \textbf{does}:
\begin{itemize}
  \item Reformulates the Collatz dynamics into a 3-adic
        admissible-inverse framework, certified line-by-line in
        Lean~4.
  \item Refutes the rigid auxiliary bound $A \ge 2q$ by explicit
        admissible word (S202).
  \item Constructs the inverse cylinder graph, the cost function
        $\Cback(m, Q)$, and a correct certificate format (S216)
        for proving lower bounds on it.
  \item Provides four explicit terminal barriers, formal in Lean,
        and three further empirical barriers awaiting routine
        scaffolding.
  \item Isolates the entire remaining analytical content as a single
        explicit conjecture (Conjecture~\ref{conj:precise}).
  \item Proves formally that this conjecture implies the slope
        barrier (Theorem~\ref{thm:slope-conditional}).
\end{itemize}

What this work \textbf{does not} do:
\begin{itemize}
  \item Prove the Collatz conjecture.
  \item Prove $\lambda_{\mathrm{asy}} < 2$ unconditionally.
  \item Resolve Conjecture~\ref{conj:precise}.
  \item Construct an analytic potential $\Phi$ closing the
        conjecture for arbitrary $(m, Q)$.
\end{itemize}

\section*{Acknowledgements}
The author thanks the Lean and Mathlib communities for the
infrastructure on which this work rests. All errors are the
author's own.

%========== References ==========
\begin{thebibliography}{99}

\bibitem{tao2022}
T.~Tao.
\newblock Almost all orbits of the Collatz map attain almost bounded
values.
\newblock {\em Forum of Mathematics, Pi}, 10(e12), 2022.

\bibitem{lagarias1985}
J.~C. Lagarias.
\newblock The 3$x$ + 1 problem and its generalizations.
\newblock {\em American Mathematical Monthly}, 92(1):3--23, 1985.

\bibitem{lagarias2010}
J.~C. Lagarias (editor).
\newblock {\em The Ultimate Challenge: The 3$x$+1 Problem}.
\newblock American Mathematical Society, 2010.

\bibitem{wirsching}
G.~J. Wirsching.
\newblock {\em The Dynamical System Generated by the 3$n$+1 Function}.
\newblock Lecture Notes in Mathematics 1681, Springer, 1998.

\bibitem{kontorovich2002}
A.~V. Kontorovich and Ya.~G. Sinai.
\newblock Structure theorem for $(d,g,h)$-maps.
\newblock {\em Bulletin of the Brazilian Mathematical Society},
33(2):213--224, 2002.

\bibitem{conway1972}
J.~H. Conway.
\newblock Unpredictable iterations.
\newblock In {\em Proceedings of the 1972 Number Theory Conference},
49--52. University of Colorado, 1972.

\bibitem{bernstein-lagarias}
D.~J. Bernstein and J.~C. Lagarias.
\newblock The 3$x$+1 conjugacy map.
\newblock {\em Canadian Journal of Mathematics}, 48(6):1154--1169, 1996.

\bibitem{cormen}
T.~H. Cormen, C.~E. Leiserson, R.~L. Rivest, and C.~Stein.
\newblock {\em Introduction to Algorithms}, 3rd edition.
\newblock MIT Press, 2009.

\bibitem{flyspeck}
T.~C. Hales \emph{et al}.
\newblock A formal proof of the Kepler conjecture.
\newblock {\em Forum of Mathematics, Pi}, 5(e2), 2017.

\bibitem{eberl-paulson}
M.~Eberl, L.~C. Paulson.
\newblock A formal proof of the prime number theorem in Isabelle/HOL.
\newblock {\em Archive of Formal Proofs}, 2018.

\bibitem{lean4}
L.~de~Moura and S.~Ullrich.
\newblock The Lean~4 theorem prover and programming language.
\newblock In {\em CADE-28}, 2021.

\bibitem{mathlib}
The mathlib Community.
\newblock The Lean Mathematical Library.
\newblock In {\em CPP 2020}, 367--381, 2020.

\bibitem{tao-blueprint}
T.~Tao \emph{et al}.
\newblock Equational theories project: a Lean blueprint methodology
for collaborative formal mathematics.
\newblock \url{https://teorth.github.io/equational_theories}, 2024--2026.

\bibitem{alphaevolve2025}
A.~Novikov \emph{et al}.
\newblock AlphaEvolve: A coding agent for scientific and algorithmic
discovery.
\newblock Google DeepMind, 2025. \url{https://deepmind.google/discover/blog/alphaevolve/}.

\bibitem{handoff}
B.~And\'ujar~Mata.
\newblock S189--S220 admissible inverse programme: handoff document.
\newblock Repository \texttt{HANDOFF\_COLLATZ\_ROADMAP.md}, May 2026.

\end{thebibliography}

\appendix

\section{Reproducibility}\label{app:reprod}

\subsection*{Public artefact}

\begin{tabular}{ll}
Repository: & \texttt{[URL to be inserted at submission time]} \\
Commit:     & \texttt{[hash to be inserted]} \\
Toolchain:  & \texttt{leanprover/lean4:v4.30.0-rc2}, Mathlib pinned in
              \texttt{lake-manifest.json}. \\
Python:     & 3.10+, standard library only. \\
\end{tabular}

\subsection*{Single-command verification}

\begin{lstlisting}
git clone <REPO URL> collatz-admissible-tree
cd collatz-admissible-tree/CollatzLean4
lake exe cache get        # downloads precompiled Mathlib
lake build                # builds the entire formal development
\end{lstlisting}

The build verifies every theorem cited above. Per-theorem axiom
inspection at the bottom of any Lean file:

\begin{lstlisting}
#print axioms CollatzLean4.Admissible.S216_barrier_for_words_outgoing
#print axioms Cert_m1_Q3_S216.barrier_m1_Q3
#print axioms Cert_m1_Q5_S216.barrier_m1_Q5
#print axioms Cert_m1_Q8_S216.barrier_m1_Q8
#print axioms Cert_m1_Q10_S216.barrier_m1_Q10
#print axioms CollatzLean4.Admissible.S202_slope_barrier_conditional
#check        CollatzLean4.Admissible.S202_one_edge_count_conjecture
\end{lstlisting}

Each \texttt{\#print axioms} call returns the empty list, modulo
Lean and Mathlib axioms.

\subsection*{Independent Python re-verification}

\begin{lstlisting}
python tools/s216_engine.py
# expected output (states):  35, 462, 24310, 200002 (or barrier
#                            certificate for higher budget), 3003.
python tools/s216_colab_debug.py   # phase-1 sanity + phase-2 frontier
\end{lstlisting}

The Phase-1 sanity script independently re-derives the state counts of
Table~\ref{tab:certs} and only proceeds to the frontier if the
sanity matches the ground truth. This is the script the author
recommends for any independent reviewer.

\subsection*{Frontier cases (\texttt{m=2, Q=9} and \texttt{m=3, Q=6})}

These are produced by \texttt{python tools/s216\_engine.py} but await
the four-step Lean scaffolding refactor described in
\texttt{CollatzLean4/AS202Lift.lean} before they enter the formal
record. The raw JSON certificates are persisted at
\texttt{tools/s216\_frontier\_results.json}.

\end{document}
